\section{방법}\label{sec:method}

본 논문의 중심 방법은 layer-adaptive K+Q steering이다. 같은 온톨로지 기저 위에서 K는 초기 레이어에만 두고 Q는 전 레이어에 유지하는 간단한 스케줄로, 정지 활성 스티어링의 반복 편향을 우회하면서 훈련 없이 추론 시 작동한다. 이 절은 먼저 기저와 세 연산자를 정의하고, 그 다음 레이어 스케줄의 직관과 비용을 서술한다.

모든 개입은 하나의 온톨로지 기저 위에서 정의된다. 도구 카탈로그의 4-facet 구조 $(\texttt{domain}, \texttt{function}, \texttt{io\_type}, \texttt{category})$를 layer-head별 basis $B \in \mathbb{R}^{d \times r}$로 정리하고, 그 사영을 $P_{\mathrm{ont}} = BB^\top$로 둔다. 이 기저는 도구를 semantic facet 축으로 분해해 attention 공간에 직접 주입할 수 있게 만든다. 본 논문은 같은 $P_{\mathrm{ont}}$ 위에서 세 연산자를 정의해 서로 비교 가능한 형태로 둔다. 정지 K-측 개입 $K' = K + \alpha K P_{\mathrm{ont}}$는 온톨로지 방향을 모든 레이어에 걸쳐 동일하게 증폭한다. signed Q-회전 $Q' = Q + \beta Q P_{\mathrm{ont}}$는 부호에 따라 focus형($\beta>0$) 또는 coverage형($\beta<0$)으로 작동한다. Layer-adaptive K+Q는 이 둘을 전 레이어에 동시에 쓰지 않고 레이어별로 다른 역할을 부여한다.

Layer-adaptive K+Q의 직관은 레이어별 민감도의 비대칭에서 나온다. 초기 레이어는 토큰 수준 의미가 형성되는 구간이라 K-측 증폭이 정답 도구 방향을 각인하는 데 유용하지만, 후반 레이어는 출력 분포가 수렴하는 구간이라 같은 K perturbation이 오히려 분포를 깨뜨린다. 따라서 본 논문은
\[
\alpha_\ell = \alpha \mathbf{1}\{\ell < \lfloor L/4\rfloor\}, \qquad \beta_\ell = \beta
\]
의 스케줄을 사용한다. K는 초기 $\lfloor L/4\rfloor$개 레이어에만 두고, Q는 전 레이어에서 유지한다. 이 스케줄은 ``초기 각인에는 K, 이후 coverage 조정에는 Q''라는 역할 분리를 명시적으로 구현한다.

본 방법의 방법론적 기여는 새로운 기저 자체가 아니라 연산자 배치에 있다. 기존 활성 스티어링은 대체로 하나의 방향을 모든 레이어에 정지적으로 주입해 왔으며, 그 결과 멀티-툴에서 반복 편향을 피하지 못했다. 본 논문은 같은 기저로도 연산자의 레이어별 배치가 결과를 결정한다는 점을 전면에 놓는다. 즉 ``더 좋은 정적 projector''가 아니라 ``멀티-툴에서 K를 쓸 수 있는 범위를 초기 레이어로 제한하고 Q와 기능을 분담시키는 operator schedule''이 본 방법의 novelty다. 이 때문에 본 방법은 SEKA류 정지 K-only 설계와 구조적으로 구분된다.

구현은 attention hook 한 번으로 충분하다. 각 레이어에서 schedule에 따라 K 또는 Q를 사영 방향으로 수정하고, 나머지 계산은 원래 모델의 forward를 그대로 사용한다. 추가 저장은 per-head basis $(L, H_{\mathrm{kv}}, d, r)$뿐이며, Qwen2.5-7B 기준 약 2.6 MB다. 가중치 업데이트와 프롬프트 확장은 필요 없다.
